\begin{abstract}
Language models can score proposed survey items cheaply, but population prompting mixes supplied group information with pretrained associations. We construct synthetic cultural agents from declared aggregate preference anchors, where a synthetic cultural agent means a group-indexed choice policy. Using Direct Preference Optimization, one adapter per profile is fit on a shared bank labeled by six GPS signs to produce a profile adapter for each group. Evaluation uses new WVS wording and a country-name-free prompt on a purposively selected development panel. We find that the trust signal transfers as a coarse sign. Agreement with human country responses remains unresolved, and a country prompt stays closer to human response distributions. The construction makes the population-specific labeling rule inspectable while leaving pretrained knowledge and synthetic-bank semantics in the policy. It provides a reusable scoring instrument whose prospective role requires independent human validation.
\end{abstract}
\noindent\textbf{Keywords:} synthetic cultural agents; aggregate anchors; direct preference optimization; survey methodology; measurement validity.
